\chapter{Chapter 7: MEC136 Master Revision Sheet}

\section{Master Formula \& Equation Reference}

\begin{table}[htbp]
\centering
\small
\begin{tabularx}{\textwidth}{l X X}
\toprule
\textbf{Formula Name} & \textbf{Mathematical Equation} & \textbf{Operational Variables \& Notes} \\
\midrule
\textbf{Representative Fraction} & $R.F. = \frac{\text{Drawing Length}}{\text{Actual Length}}$ & Units of numerator and denominator must match before reducing fraction. \\
\addlinespace
\textbf{Length of Scale} & $L.O.S. = R.F. \times \text{Max Length}$ & Used to calculate physical line length for scale on paper. \\
\addlinespace
\textbf{Lettering Spacing} & Spacing $= \frac{1}{5} h, \quad \text{Line} = 0.5h \text{ to } 1h$ & Single-stroke vertical Gothic, $7:4$ height-to-width ratio. \\
\addlinespace
\textbf{Isometric Scale Factor} & $\frac{\text{Iso Length}}{\text{True Length}} = \frac{\cos 45^\circ}{\cos 30^\circ} = \sqrt{\frac{2}{3}} \approx 0.816$ & Isometric Length $= 0.816 \times \text{True Length}$. \\
\addlinespace
\textbf{Front View Apparent Length} & $a'b' = TL \cos \phi$ & Where $\phi$ is inclination of line to VP. \\
\addlinespace
\textbf{Top View Apparent Length} & $ab = TL \cos \theta$ & Where $\theta$ is inclination of line to HP. \\
\addlinespace
\textbf{Cylinder Stretch-out Length} & $L = \pi D$ & Parallel Line Method base perimeter for cylinder. \\
\addlinespace
\textbf{Cone Sector Angle} & $\theta = \left( \frac{r}{L} \right) \times 360^\circ$ & Where $r = \text{base radius}$, $L = \text{true slant height}$. \\
\bottomrule
\end{tabularx}
\caption{Master Formula Reference Sheet.}
\label{tab:master-formulas}
\end{table}

---

\section{Projection Rules \& Quadrant Quick Reference}

\begin{table}[htbp]
\centering
\small
\begin{tabularx}{\textwidth}{l p{4.5cm} X}
\toprule
\textbf{System / Quadrant} & \textbf{View Layout Rules} & \textbf{Key Characteristics} \\
\midrule
\textbf{First Angle (ISO)} & FV above $XY$, TV below FV, LSV to right of FV & Object between Observer and Plane. Standard in India/Europe. \\
\addlinespace
\textbf{Third Angle (ANSI)} & FV below $XY$, TV above FV, RSV to right of FV & Plane between Observer and Object. Standard in USA/Canada. \\
\addlinespace
\textbf{$1^{\text{st}}$ Quadrant Point} & $a'$ above $XY$, $a$ below $XY$ & Point above HP, in front of VP. \\
\addlinespace
\textbf{$2^{\text{nd}}$ Quadrant Point} & $a'$ above $XY$, $a$ above $XY$ & Views overlap above reference line $XY$. \\
\addlinespace
\textbf{$3^{\text{rd}}$ Quadrant Point} & $a'$ below $XY$, $a$ above $XY$ & Point below HP, behind VP. \\
\addlinespace
\textbf{$4^{\text{th}}$ Quadrant Point} & $a'$ below $XY$, $a$ below $XY$ & Views overlap below reference line $XY$. \\
\bottomrule
\end{tabularx}
\caption{Projection System and Point Location Rules.}
\label{tab:projection-rules-master}
\end{table}

---

\section{AutoCAD Command Cheat Sheet}

\begin{table}[htbp]
\centering
\small
\begin{tabularx}{\textwidth}{c c X X}
\toprule
\textbf{Command} & \textbf{Shortcut} & \textbf{Primary Purpose} & \textbf{Key Option / Trap} \\
\midrule
\texttt{UNITS} & \texttt{UN} & Set precision and unit types & Set Insertion Scale to Millimeters. \\
\texttt{LIMITS} & — & Set drawing boundaries & Must follow with \texttt{ZOOM ALL} (\texttt{Z} $\to$ \texttt{A}). \\
\texttt{OSNAP} & \texttt{F3} & Object Snap toggles & Enable Endpoint, Midpoint, Center. \\
\texttt{ORTHO} & \texttt{F8} & Restrict motion to horizontal/vertical & Essential for straight projection lines. \\
\texttt{LINE} & \texttt{L} & Create straight line segments & Polar syntax: \texttt{@length<angle}. \\
\texttt{PLINE} & \texttt{PL} & Create connected 2D polyline & Forms a single unified object. \\
\texttt{TRIM} & \texttt{TR} & Cut objects to boundary edge & Select cutting edge $\to$ click object. \\
\texttt{FILLET} & \texttt{F} & Round sharp corner edges & Set radius first: \texttt{F} $\to$ \texttt{R} $\to$ Value. \\
\texttt{CHAMFER} & \texttt{CHA} & Bevel sharp corner edges & Set distance first: \texttt{CHA} $\to$ \texttt{D} $\to$ Values. \\
\texttt{STRETCH} & \texttt{S} & Elongate geometry & MUST use Green Crossing Window. \\
\texttt{HATCH} & \texttt{H} & Fill section with pattern & Pattern \texttt{ANSI31} at $45^\circ$, Pick Points. \\
\texttt{EXTRUDE} & \texttt{EXT} & Convert 2D profile to 3D solid & Input must be closed polyline/region. \\
\texttt{REVOLVE} & \texttt{REV} & Sweep 2D profile around axis & Requires closed profile and axis line. \\
\texttt{PRESSPULL} & — & Dynamic 3D volume pull & Auto-detects enclosed coplanar area. \\
\texttt{SUBTRACT} & \texttt{SU} & Cut 3D volume & Select KEEP solid $\to$ Select REMOVE solid. \\
\texttt{VIEWBASE} & — & Generate 2D views from 3D model & Auto-links Model Space to Layout Tab. \\
\bottomrule
\end{tabularx}
\caption{AutoCAD Command Cheat Sheet.}
\label{tab:autocad-cheat-sheet}
\end{table}

---

\section{Common Exam Traps \& Pitfalls}
\begin{enumerate}
    \item \textbf{IC 7402 / Sectioning Misconceptions:}  
    In Sectional Views, thin webs/ribs cut longitudinally are **NOT** hatched. Fasteners (bolts, nuts, screws) and solid shafts cut longitudinally are **NOT** hatched.
    \item **Isometric View vs. Projection:**  
    Isometric **View** uses True Scale ($1:1$). Isometric **Projection** uses reduced scale factor ($0.816$). Do not confuse the two when reading exam instructions!
    \item **Line Projection Projector Alignment:**  
    In the Rotating Line Method, Front View end point $b'$ and Top View end point $b$ **MUST lie on the same vertical projector line**. If they do not align, the construction is incorrect.
    \item **AutoCAD Stretch Command Error:**  
    Selecting objects with a single click or a Blue Window fails for \texttt{STRETCH}. You **MUST** use a Green Crossing Window (dragged right-to-left).
\end{enumerate}
